
\chapter{Управление отказами}\label{ch:decline-management}

Когда покупатель видит отказ на странице оплаты, для него всё выглядит
просто: платёж не прошёл. Для платёжной команды на этом простота
заканчивается. Один отказ означает \enquote{подожди и~попробуй ещё раз},
другой -- \enquote{сломался маршрут, а~не счёт клиента}. Задача --
различить эти случаи до того, как система начнёт бессмысленно жечь
попытки и раздражать держателя карты.

\section{Анатомия отказа}\label{sec:decline-anatomy}

На кассе, в~приложении и в~дашборде все отказы выглядят одинаково:
операция не состоялась. Но дихотомия \enquote{мягкий /
жёсткий отказ} остаётся эвристикой продавца. Visa делит
коды отказов (decline response codes) на категории
\enquote{Issuer will never approve}, \enquote{Issuer cannot approve at this
time}, \enquote{Data quality} и \enquote{Generic}; Mastercard
отдельно запрещает дополнительные CNP-авторизационные запросы для части
кодов отказов~\cite{visa-core-rules,mc-tpr}. Практический вопрос
один: какое действие разумно после конкретного ответа.

Мягкий отказ (soft decline) означает, что транзакция отклонена по
временн\'{о}й или устранимой причине, временный лимит, поведенческий
фильтр, недоступность сервиса. Эмитент как будто говорит:
\enquote{не сейчас, попробуй ещё позже}. К~этой группе на практике чаще
всего относят \texttt{51}: Insufficient Funds, \texttt{61}:
Exceeds Withdrawal
Amount Limit, \texttt{65}: Exceeds Withdrawal
Frequency Limit и~часть общих (\textit{generic}) кодов вроде \texttt{05}, если схема вообще
допускает повтор и~ТСП остаётся в~пределах лимитов
повторов~\cite{visa-core-rules}.

Жёсткий отказ (hard decline) в~такой операционной модели означает, что ту
же операцию с~теми же данными повторять бессмысленно или прямо запрещено
правилами. Надёжные примеры здесь: \texttt{04}: Pick Up, \texttt{14}:
Invalid Account Number\footnote{В~актуальной редакции Visa Core Rules
код~\texttt{14} отнесён к~категории~1 (Issuer will never approve);
операционно он эквивалентен жёсткому отказу -- повтор с~теми же
реквизитами бессмысленен, нужен новый PAN.},
\texttt{41}: Lost Card и \texttt{43}: Stolen Card.
Они либо указывают на неверные реквизиты, либо на карту, с~которой
продолжать операцию нельзя. Здесь нужен другой платёжный реквизит или
ручная проверка~\cite{visa-core-rules,mc-tpr}.

Есть и~коды, которые ломают простое деление. \texttt{54} Expired Card для
Visa относится к~категории Data quality: схема прямо трактует
его как случай, где нужно перепроверить платёжную информацию. У
Mastercard правило сформулировано иначе: для CNP ТСП не
должен инициировать дополнительные авторизационные запросы по той же
операции с~тем же PAN и~сроком действия -- повторять тот же запрос без новых реквизитов
нельзя~\cite{visa-core-rules,mc-tpr}.

Отдельно стоят технические и маршрутные коды. \texttt{91}: Issuer
Unavailable и \texttt{96}: System Malfunction говорят скорее о~доступности
эмитента или маршрута, чем о~состоянии счёта. Для них схема допускает
повторные попытки в~пределах лимитов, но смысл работы здесь другой:
сначала убедиться, что проблема действительно в~доступности пути, а~не в
самой карте. \texttt{01}: Refer to Card Issuer тоже не приглашение к
слепому повтору: обычно это знак, что нужен иной канал или
дополнительное действие со стороны держателя либо
эмитента~\cite{visa-core-rules}.

\begin{table}[tbp]
  \begingroup
  \scriptsize
  \setlength{\tabcolsep}{4pt}
  \renewcommand{\arraystretch}{1.2}
  \makebox[\textwidth][l]{%
    \begin{minipage}{\fullwidth}
      \begin{tabularx}{\fullwidth}{@{}Y Y Y@{}}
        \toprule
        \cellcolor{Bad!16}\textbf{Не повторять ту же операцию}
        &
        \cellcolor{Warn!18}\textbf{Нужны новые данные / пауза}
        &
        \cellcolor{AccentA!14}\textbf{Технический / другой путь}
        \\
        \midrule
        \cellcolor{Bad!4}\textbf{04 Pick Up}\newline stop; запросить другой credential
        \newline\textbf{14 Invalid Account}\newline проверить PAN
        \newline\textbf{41 Lost Card}\newline карта утеряна; stop
        \newline\textbf{43 Stolen Card}\newline карта украдена; stop
        &
        \cellcolor{Warn!5}\textbf{54 Expired Card}\newline обновить expiry
        / PAN; в Mastercard CNP не повторять тот же PAN+expiry
        \newline\textbf{51 Insufficient Funds}\newline пауза; не делать blind loop
        \newline\textbf{61 Exceeds Limit}\newline проверить сумму и~лимит
        \newline\textbf{62 Restricted Card}\newline смотреть контекст и~правила схемы
        \newline\textbf{65 Freq. Limit}\newline дождаться сброса лимита
        \newline\textbf{05 Do Not Honor}\newline общий отказ; только в
        пределах caps &
        \cellcolor{AccentA!5}\textbf{91 Issuer Unavailable}\newline
        проблема доступности route / issuer
        \newline\textbf{96 System Malfunction}\newline технический сбой;
        сначала final state
        \newline\textbf{01 Refer}\newline нужен иной канал или действие банка
        \newline\textbf{timeout}\newline подтвердить финальный статус,
        затем controlled retry \\
        \bottomrule
      \end{tabularx}
      \par\smallskip
      \fcolorbox{Ink!18}{Soft}{%
        \parbox{\dimexpr\fullwidth-2\fboxsep-2\fboxrule\relax}{%
          \textcolor{Muted}{Точная схема сложнее этой таблицы: Visa делит
            коды отказов на Category~1/2/3/4,
            а Mastercard отдельно запрещает часть CNP retries. Здесь
            показан только первый операционный вопрос:
            прекращать операцию, обновлять данные или работать с~доступностью
          маршрута.}%
        }%
      }
    \end{minipage}%
  }
  \endgroup
  \caption{Рабочая классификация кодов ответа по следующему действию.}\label{fig:decline-codes}
\end{table}

Сначала определяется обратимость
отказа, потом решается вопрос о~повторе. Код ответа --
сжатая операционная подсказка,
которая всегда сверяется с~правилами конкретной схемы.

\section{Что на самом деле означают коды ответа}\label{sec:decline-response-codes}

Код ответа -- ответ эмитента, но не его объяснение. У~платёжной
команды всегда есть соблазн прочитать его как точную причину отказа.
Почти всегда это ошибка. За одним и~тем же кодом стоят десятки причин,
и~без дополнительных данных истинную причину не узнать.

\texttt{05} Do Not Honor: самый распространённый и самый неинформативный
код. Формально он означает \enquote{не проводите операцию}, но почти ничего не
объясняет. За ним скрываются антифрод, дневной лимит, несоответствие
поведенческому профилю клиента, временная блокировка в~приложении или
внутренняя политика эмитента по конкретному MCC. Часть причин
временная и исчезает сама, часть постоянная; поэтому
\texttt{05} трактуют с~осторожностью. У Visa общие коды отказов (generic decline
response codes) сведены в~отдельную Category~4 и~должны использоваться
только там, где не подходит более точное значение. Этого достаточно,
чтобы трактовать \texttt{05} как отказ без расшифровки~\cite{visa-core-rules}.
Эмитент намеренно не раскрывает причину через конкретный код: детальный
код -- информация, полезная атакующему. Если \texttt{51} точно
означает нехватку средств, злоумышленник с~украденными картами отбросит
\enquote{пустые} и сосредоточится на \enquote{богатых}. Возвращая \texttt{05},
эмитент лишает атаку перебором краденых учётных данных (\textit{credential stuffing}) сигнала о~состоянии счёта.

\texttt{51} Insufficient Funds тоже не сводится к~буквальному бытовому
смыслу. Для дебетовых карт он часто означает нехватку
средств на счёте. Для кредитных карт это превышение лимита,
а~у~некоторых эмитентов -- универсальная форма мягкого
отказа, когда раскрывать точную причину они не хотят.

\texttt{54} Expired Card как раз показывает, почему буквальное чтение опасно.
В~таблице категорий кодов отказа Visa Core Rules это код класса
Data quality, а~в Visa Token Service
PAN Lifecycle прямо предусмотрены обновление PAN, срока действия и VAU
updates. Поэтому в потоках с~повторными списаниями и~сохранёнными реквизитами карт правильный вопрос
после \texttt{54} часто не \enquote{послать ли ещё раз}, а
\enquote{получили ли мы уже обновлённые реквизиты}. У Mastercard ту же
операционную задачу закрывает Automatic Billing Updater, который
публично позиционируется как способ держать платёжные реквизиты
актуальными
и избегать
отказов~\cite{visa-core-rules,visa-vts,visa-vau,mc-bill-pay-solutions}.

\texttt{91} и \texttt{96}: технические коды, которые ничего не говорят о~состоянии
карты или счёта. Они означают, что эмитент или маршрут недоступны;
транзакцию уместно рассматривать как вопрос доступности, а~не
кредитного решения. В Visa оба кода попадают в~классы, где
повторы разрешены в~пределах схемных лимитов. Если такие коды
доминируют в~статистике, проблема лежит в~сети,
маршрутизации или доступности хоста. Карта здесь вообще ни при чём~\cite{visa-core-rules}.

Отдельная категория: мягкие отказы, связанные с~аутентификацией.
Код~\texttt{65} в~исходной ISO~8583/Visa-семантике
означает Exceeds Withdrawal Frequency Limit (нужно дождаться
сброса лимита частоты); в~европейском SCA-контексте Mastercard
тот же \texttt{65} используется в~качестве soft decline,
по которому эмитент готов одобрить операцию после повторной
аутентификации. Поток после такого отказа: шлюз получает \texttt{65},
распознаёт его как retryable-with-3DS, инициирует
аутентификацию (AReq~$\to$~DS~$\to$~ACS), получает результат
(ECI + CAVV) и повторяет авторизационный запрос с~данными
аутентификации. Эмитент видит тот же PAN, ту же сумму, но
с~подтверждением SCA -- и~одобряет. Для продавца разница
между слепым retry того же запроса и retry с 3DS -- это
разница между бесконечным циклом отказов и восстановленной
конверсией~\cite{mc-tpr}.

\practicenote{%
  Хороший мониторинг отказов сегментирует коды ответа не по двум
  группам, а~хотя бы по трём-четырём: жёсткие, мягкие финансовые,
  мягкие поведенческие и технические. Сводный дашборд с~одним числом
  доли отказов скрывает структуру: 15\% отказов, половина из которых
  приходится на технические коды, означают совсем другую задачу, чем
  те же 15\% при преобладании необратимых отказов.

}
\section{Повторы после отказа: время, частота, каналы}\label{sec:decline-retry}

На повторах расходятся интересы продавца и платёжной схемы. Продавец
хочет вернуть конверсию; схема хочет, чтобы сеть не превращалась
в~генератор бессмысленных попыток. Ошибка в~логике повторов
обходится потерянной выручкой и~штрафами мониторинга.

Каждый бессмысленный повтор
потребляет авторизационную ёмкость схемы; при миллионах подписочных
ТСП даже два лишних запроса в~месяц на клиента создают значимую
нагрузку. Схема ставит барьер там, где маргинальная польза повтора
уже близка к~нулю -- это равновесие между конверсией продавца
и стабильностью сети.

Правила схем задают прежде всего внешние границы:
какие коды вообще можно повторять, а~какие нельзя. Visa Core Rules
сводят коды отказов в~таблицу категорий с~разными лимитами
повторов для ТСП; Mastercard для
CNP отдельно запрещает дополнительные авторизационные запросы по той же
транзакции после ряда кодов отказа, включая \texttt{04},
\texttt{14}, \texttt{41},
\texttt{43} и \texttt{54}~\cite{visa-core-rules,mc-tpr}.

Схемы почти не говорят ТСП
\enquote{повтори через 17 минут}; они говорят \enquote{не делай слепой
цикл повторов и~не выходи за лимиты}. Если отказ связан с~нехваткой
средств или лимитом, серия одинаковых запросов не создаёт деньги и~не
расширяет лимит. Если код технический, бессмысленная очередь повторов
лишь засоряет наблюдаемость и повышает риск мониторинга.

Логика повторов сводится к~двум вопросам. Допускает ли схема повтор
этого кода -- или нужны новые данные (срок действия, токен, 3DS)?
Если повтор допустим, какой интервал и~какой потолок безопасны для
конкретного продавца, PSP и вертикали? Mastercard на
публичных продуктовых страницах продаёт это как intelligent retry
strategies~\cite{mc-bill-pay-solutions}.

На практике коды ответа группируются в~пять классов
по допустимости повтора:

\begin{itemize}
  \item \textbf{Жёсткий отказ} (\texttt{04}, \texttt{14},
    \texttt{41}, \texttt{43}): карта недействительна, изъята
    или украдена. Повтор запрещён; нужен новый платёжный
    инструмент.
  \item \textbf{Мягкий финансовый} (\texttt{51}): нехватка
    средств. Повтор допустим, но не ранее чем через 24--48~часов
    (когда баланс может измениться); цикл с~интервалом в~минуты
    бесполезен.
  \item \textbf{Мягкий аутентификационный} (\texttt{65},
    \texttt{1A}): требуется SCA. Повтор через 3DS, а~не повтор
    того же запроса.
  \item \textbf{Технический} (\texttt{91}, \texttt{96}):
    эмитент или маршрут недоступен. Повтор допустим через
    15--60~минут; каскад через другого эквайера эффективен,
    только если альтернативный маршрут ведёт к~другому хосту
    эмитента.
  \item \textbf{Непрозрачный} (\texttt{05}): причина неизвестна.
    Серия идентичных запросов повышает риск попадания
    в~мониторинг схемы.
\end{itemize}

Канал повтора тоже меняет исход. Visa и Mastercard описывают
токенизацию и оптимизацию одобрений как способы дать эмитенту более
чистый сигнал и~убрать лишнюю фрикцию. Повтор через
сетевой токен или через более удачный путь маршрутизации
не равен простому повтору
PAN-запроса~\cite{visa-vts,visa-tokenization-security,mc-acquirers-psps,mc-pop}.

\begin{figure}[tbp]
  \centering
  \includefiguresvg[width=\linewidth]{assets/figures/ch25-decline-management/decline-retry-tree}
  \caption{Рабочая схема действий после отказа в~авторизации.}\label{fig:decline-retry-tree}
\end{figure}

Хорошая логика повторов уменьшает количество бессмысленных попыток.
Снижение самих отказов -- отдельная задача. Смешивать эти два
направления опасно.

\section{Сетевой токен и PAN: дополнительный контекст доверия}\label{sec:decline-tokens}

Если продавец работает через сетевой токен, а~не хранит PAN,
эмитент получает более управляемый контекст. Visa Token Service
описывает токен как отдельный цифровой идентификатор, с
которым эмитент задаёт token variables, определяет авторизованных
token requestors и запрашивает token details, включая device и risk
information. Токен может быть
ограничен конкретным мобильным устройством, \textit{e-commerce}-ТСП или
количеством покупок~\cite{visa-vts,emvco-payment-tokenisation-guide}.

Сетевой токен несёт управляемый контекст о~том, где и~как реквизит
может использоваться. Visa в~публичных security materials
исторически озвучивала рост доли одобрений (approval rate) около
2,5\,\% на горизонте 2020--2022~годов; в~обновлённых материалах
2024--2026~годов диапазон, на который ссылается Visa, доходит
до 4--6\,\%. Mastercard на странице для acquirers и PSP пишет
о 3--6~п.п. роста доли одобрений у tokenized digital
payments~\cite{visa-tokenization-security,mc-acquirers-psps}. Цифры
принадлежат самим схемам и~не обещают магии каждому продавцу, но
общий вывод устойчив: токен даёт эмитенту более чистый сигнал,
чем обычный PAN.

Самый приземлённый выигрыш виден в~жизненном цикле реквизитов. Visa PAN
Lifecycle обновляет PAN, срок действия и выполняет VAU
updates. У Mastercard Automatic Billing Updater подаётся как способ
держать платёжные реквизиты актуальными и избегать отказов.
\texttt{54} Expired Card в потоках с~сохранёнными реквизитами карт чаще означает не
\enquote{карта умерла}, а \enquote{у~нас устарела ссылка на
реквизит}~\cite{visa-vts,visa-vau,mc-bill-pay-solutions}.

Миграция на токены откладывается по скучным, но устойчивым причинам.
Нужны интеграция и сертификация; если текущий поток кое-как держится,
миграцию откладывают до заметной боли в~виде доли отказов и~оттока;
многие продавцы не управляют токенами напрямую и
зависят от того, как быстро двигается их PSP. Токены остаются
вопросом и безопасности, и операционной зрелости.

Подробнее об архитектуре сетевых токенов в~главе~\ref{ch:tokenization}.

\section{Маршрутизация через эквайера как инструмент конверсии}\label{sec:decline-routing}

Маршрутизация через эквайера имеет смысл лишь потому, что приём
платежей -- не однотрубная система. В~правилах Mastercard есть
отдельные разделы про authorization routing и domestic transaction
routing, а~в~публичных материалах Mastercard для acceptance community
routing и payment optimization вынесены в~самостоятельный слой
продуктов~\cite{mc-tpr,mc-acquirers-psps,mc-pop}.
Путь до эмитента влияет на исход и~не сводится к~вопросу
\enquote{какой PSP мы подключили}.

У~разных эквайеров различаются локальное покрытие, подключения, качество
формирования сообщения и сервисы оптимизации. Сетевая оптимизация
одобрений работает именно здесь: более удачный маршрут авторизации или
лучше сформированное сообщение~\cite{mc-pop}.

Страна BIN и~страна ТСП определяют, транзакция внутренняя (\textit{domestic})
или трансграничная (\textit{cross-border}): внутренний маршрут через локального
эквайера даёт эмитенту знакомый контекст и~часто лучше конвертирует.
MCC влияет на доступные программы \textit{interchange} и на набор
антифрод-правил эмитента. Валюта транзакции и валюта расчёта
важны, если платформа поддерживает DCC. Историческая доля одобрений
на конкретном эквайер--BIN--MCC сочетании даёт основу для
статистически обоснованного решения о~выборе маршрута в~реальном
времени.

Принцип выбора маршрута -- функция ожидаемого дохода:
$P(\text{approve} \mid \text{route}) \times \text{amount} -
\text{cost}(\text{route})$. Маршрут с~более высокой комиссией
оправдан только при достаточно высоком $P(\text{approve})$. В~зрелых
системах -- ML-модель, переобучаемая на исторических авторизациях;
в~простых -- скоринговая (\textit{score}) таблица по BIN~$\times$~MCC~$\times$~страна.

Различие cascade routing и intelligent routing как двух паттернов
маршрутизации, корректная обработка задержавшегося ответа первого
эквайера и~ограничения каскада по типам отказов разобраны
в~\ref{sec:gw-routing}. Здесь существенно операционное следствие:
если все эквайеры в~каскаде вернули отказ, транзакция завершается
с~последним полученным кодом, и~повторно прогонять её через тех же
эквайеров запрещено правилами повторов схем. Правильное действие --
показать ошибку или предложить другой платёжный метод.

Свобода маршрутизации не безгранична. Она остаётся внутри
правил схем, локальной маршрутизации и interchange; она -- не лицензия
на маскировку географии ТСП или искусственное
конструирование тарифного режима. Легитимная цель routing --
повысить долю одобрений и сократить операционные
потери~\cite{visa-core-rules,mc-rules,mc-tpr}.

Smart routing -- слой дисциплины поверх платёжной
инфраструктуры: сначала ограничения пути, затем лучший доступный маршрут.

\section{Доля одобрений по вертикалям}\label{sec:decline-benchmarks}

85\% одобрений: это хорошо или тревожно? Ответ зависит от
вертикали, географии, типа карты и~самой модели оплаты.

Сначала договариваются о~знаменателе и о~самом разрезе. Сетевые
правила и мониторинги сами считают approval performance не одним
универсальным числом: Mastercard отслеживает долю одобрений
в CNP- и трансграничных контекстах через Merchant Advice
Code (MAC), Authorization Optimizer и связанные
мониторинговые программы. Это предупреждает против сравнения
всех потоков одной агрегированной метрикой~\cite{mc-tpr}.

Сами диапазоны approval rate стареют быстрее, чем успевает
выйти книга: они зависят от вертикали, страны, доли 3DS,
качества токенизации, текущих волн мошенничества и даже
календарного окна. По этой же причине универсальные ориентиры
\enquote{нормальной} доли одобрений по вертикалям в~книге
не приводятся: даже опубликованные числа быстро устаревают
и не учитывают разрезов конкретного продавца. Конкретные
проценты для своего потока сверяются по данным PSP или
эквайера. Системные просадки относительно собственной
исторической нормы по вертикали -- сильный сигнал: неправильный
retry, отсутствие 3DS на мягкий отказ, невыверенный BIN-маршрут
или устаревшие реквизиты в \textit{card-on-file}.

Сравнивать розницу с~сервисами путешествий, подписки с~разовыми
интернет-платежами, а внутренний поток с~трансграничным в~одной колонке
бессмысленно. В recurring добавляется контекст
сохранённого реквизита: если цепочка согласия не оформлена по
правилам SCF, эмитент
видит обычный CNP-запрос, а~не продолжение уже данного согласия~\cite{visa-scf,mc-scf}.

\importantnote{%
  Долю одобрений нельзя читать одним числом по продукту. Метод только
  один: смотреть разрезы по странам, вертикалям и типам интеграции,
  иначе реальные причины отказов прячутся внутри агрегата.

}
\practicenote{%
  Измеряйте долю одобрений раздельно по сегментам: новые и вернувшиеся
  клиенты, карты с 3DS и~без неё, PAN и~токен, внутренний рынок и
  трансграничные операции. Агрегированная метрика почти бесполезна для
  оптимизации. Сегментация быстро показывает, где сосредоточены
  отказы: как правило, основную массу проблем даёт не весь поток, а
  сравнительно узкий класс транзакций с~особыми характеристиками.

}
